\section{Related Work}
\label{sec:related_work}

%---------------------------------------------------------------------
\begin{figure*}[ht]
  \centering
  \includegraphics[width=\linewidth]{./figs/overview.pdf}\\
  \caption{\label{overview} 
    \textbf{Overview of the LaGen framework.} It first obtains estimates of the current frame's foreground and background point clouds via the Scene Decoupling Estimation (SDE) module. Subsequently, all three-dimensional information is projected to the two-dimensional range image space via spherical mapping. These range-view representations are then processed by a pretrained VAE and a noise modulation (NM) module, before being passed to the U-net network to generate the LiDAR scene for the current frame. Afterward, the current generated scene is iteratively used as the input to the model, thereby completing the autoregressive process.}
  \vspace{-5mm}  
  \end{figure*}
%---------------------------------------------------------------------

\paragraph{Generative Models.}
Generative models aim to learn the underlying data distribution and generate novel samples. 
Generative Adversarial Networks (GANs) \cite{goodfellow2014generative} are a prominent example. 
GANs alternately train a pair of competing networks, the generator and the discriminator, to minimize the adversarial objective.
%Although GANs achieve high data synthesis quality, their training process is often unstable.
However, GANs often suffer from unstable training dynamics.
Diffusion models (DMs) \cite{sohl2015deep} have shown outstanding performance in the field of image synthesis. 
Mainstream diffusion models include score-based models \cite{song2019generative, song2020improved, song2020score} and Denoising Diffusion Probabilistic Models (DDPM) \cite{ho2020denoising, nichol2021improved, kingma2021variational, saharia2022photorealistic}.
They enable stable training using simple objective functions, addressing challenges associated with GANs.
Recently, Latent Diffusion Models (LDMs) \cite{rombach2022high} have been proposed. 
%LDMs construct a perceptually equivalent latent space using a pretrained autoencoder and perform the diffusion process within this latent space, greatly improving generation efficiency.
They perform the diffusion process in the latent space, greatly improving generation efficiency.
They have also enabled significant advances in multimodal conditional generation \cite{nichol2021glide, ramesh2022hierarchical, zhang2023adding}.

\vspace{-4mm}

\paragraph{LiDAR Data Generation.}
Generating LiDAR point clouds is a three-dimensional data generation task. 
Caccia et al. \cite{caccia2019deep} were the first to represent LiDAR data as range images and employed Variational Autoencoders (VAEs) \cite{kingma2013auto} and Generative Adversarial Networks (GANs) \cite{ goodfellow2014generative} for generation.
Subsequently, LiDARGen \cite{LiDARGen2022} introduced a score-based LiDAR data generation model. 
UltraLiDAR \cite{xiong2023ultralidar} voxelizes LiDAR point clouds and employs VQ-VAE \cite{yu2021vector} for generation.
LiDM \cite{ran2024towards} separately generates scenes, objects, and trajectories, and simulates ray casting to produce point clouds.
R2DM \cite{nakashima2024lidar}, based on Denoising Diffusion Probabilistic Models (DDPMs) \cite{song2020denoising}, further improves the fidelity of generated point clouds.
%OLiDM \cite{yan2025olidm} proposes two modules: the Object-Scene Progressive Generation module and the Object Semantic Alignment module, to enhance model performance on downstream perception tasks. 
RangeLDM \cite{hu2024rangeldm} introduces a range-guided discriminator, enabling fast generation through a latent diffusion model.
However, most of these frameworks are only capable of generating isolated single-frame LiDAR data, or lack real-time interactivity.

\vspace{-3mm}

\paragraph{LiDAR Data Forecasting.}
Point cloud prediction \cite{mersch2022self, weng2021inverting, weng2022s2net, wilson2023argoverse} is a fundamental self-supervised task in autonomous driving, where multiple frames of historical point clouds are used as input to predict future point clouds. 
Early works focused on estimating scene flow between consecutive point clouds to forecast future point cloud sequences \cite{liu2019flownet3d, zhai2021optical}. 
However, such approaches typically incur high computational costs, limiting their practical applicability.
Recent studies have started to directly predict from raw point cloud data \cite{luo2023pcpnet, pal2024atppnet, weng2021inverting, weng2022s2net}. 
4D-Occ \cite{khurana2023point} predicts point clouds from the perspective of geometric occupancy prediction, which further improves prediction accuracy.
ViDAR \cite{yang2024visual} focuses on visual point cloud prediction, using historical images to predict future point clouds, and thereby improves performance in downstream applications to some extent.
These prediction-based frameworks typically require multiple frames of historical input, yet still cannot accurately predict long-horizon LiDAR scenes and lack interactivity.